\documentclass[openany,fontset=fandol]{ctexbook}
\usepackage{amsmath,amssymb,amsthm,enumitem}
\usepackage[unicode,hidelinks]{hyperref}
\setlength{\emergencystretch}{2em}
\raggedbottom
\DeclareMathOperator{\Spec}{Spec}
\DeclareMathOperator{\Supp}{Supp}
\newcommand{\FF}{\mathcal F}
\newcommand{\GG}{\mathcal G}
\newcommand{\calP}{\mathcal P}
\begin{document}
\title{代数几何基础第五卷预备笔记\\法文覆盖审计：超平面截面与多截面构造}
\author{第97批\quad 法、英、中文三方校订}
\date{2026年11月24日}
\maketitle

\section*{范围与接续}
本批核对法文物理26、27页，及英编稿物理18页底至19页的相应段。正文译校覆盖7.1从节首至有限域换底及引出例子的过渡句。两处划销旧稿另列；物理28页起关于4.2、4.3的现行补记和命题7.2留待下一批。第三辑1158页合订本未改动。

\section*{7\quad 应用于构造指定类型的超平面截面或多截面}

\paragraph{7.1.}
设 $S=\Spec k$，且 $k$ 是无限域。那么 $P^\vee$ 的任一非空开子集都含有 $k$-有理点。采用4.1的记号，若由可构造性质 $\calP$ 定出的 $E$ 含泛点 $\eta$，则 $E$ 含一个 $k$-有理点。因此，$X$ 存在一个在 $k$ 上定义且满足 $\calP$ 的超平面截面。

现在让 $S$ 仍为任意基底。给定 $s\in S$ 及纤维 $P^\vee_s$ 的一个 $k(s)$-有理点 $\xi_0$，可以在 $s$ 的一个开邻域 $U$ 上找到 $P^\vee\to S$ 的截面 $\xi$，使它经过 $\xi_0$。仍如4.1以可构造性质 $\calP$ 定义 $E$。若已知 $E$ 是开的（例如由第5节得到这一点），且 $\xi_0\in E$，则充分缩小 $U$ 后，$\xi$ 的像落在 $E$ 中。于是得到 $X$ 在 $U$ 上的一个超平面截面 $Y_\xi$，使每个 $t\in U$ 上的纤维 $Y_{\xi(t)}$ 都满足 $\calP$。\footnote{法26将初始点手改成 $\xi_0$、将开放性的依据改指第5节，并补记 $Y_\xi$；法27首行手补 $Y_{\xi(t)}$。区别参数丛的截面 $\xi$ 和 $X$ 的超平面截面 $Y_\xi$。}

若事先没有指定 $\xi_0$，而 $k(s)$ 是无限域，可把前两点结合起来，在 $s$ 的某个开邻域上构造具有上述性质的超平面截面。\footnote{这里仍须在 $P^\vee_s$ 中有一个非空的好参数开集，例如纤维内相应好参数集可构造且含其泛点。单凭 $k(s)$ 无限，不能推出有好截面。}
最后，利用第5节可得到如下判据：在 $X$（相应地 $\FF$）相对 $S$ 平坦并满足该节判据的其余条件时，$Y_\xi$（相应地 $\GG_\xi$）也相对 $S$ 平坦。因而可用 $Y_\xi$ 代替 $X$，重复上述构造。比如，在一定条件下，可以逐步构造 $X$ 在给定点 $s$ 的某个开邻域 $U$ 上的一条“多截面” $S'$，使 $S'\to U$ 有限、平坦且各纤维满足 $\calP$。\footnote{“逐步构造”取法文清楚可辨的 \emph{de proche en proche}，不将英译问号处译作另有极限的“逐次逼近”。法稿说“例如，在一定条件下”，未在这里给出无条件多截面定理。}

若 $k(s)$ 是有限域，可能须先作满射的 étale 换底 $S'\to U$，再作上述构造。确实，与7.1开头相同，考虑 $P^\vee$ 中的非空开集；若 $k$ 有限，该开集未必有 $k$-有理点，但必有闭点，给出取值于 $k$ 的有限扩域 $k'$ 的点。这个扩域必可分。令 $k=k(s)$，在 $s$ 的某个开邻域 $U$ 上作与剩余域扩张 $k'/k(s)$ 对应的适当有限 étale 覆盖 $S'\to U$，使
\[
 S'_s\simeq\Spec k'.
\]
于是经换底 $S'\to S$，可在 $s$ 上方唯一的点 $s'\in S'$ 处归结到有一个有理好参数点的情形。\footnote{法27、英19均写“6.1开头”，实际说的是本条7.1起首的非空开集与有理点。该段把参数空间的非空开集和基底邻域都写成 $U$，译文分开。有限扩域并不会变成无限域；有利之处是选定的好闭点在新剩余域上成为有理点。这里的 $S'$ 表示换底覆盖，不是上一段的目标多截面。}
下面给出用这些方法得到的一些命题的例子。

\section*{译校补证}
\paragraph{无限域与局部提升.}
$P^\vee$ 的标准仿射图上，非空主开集由一个非零多项式 $f$ 不消失来刻画。无限域上这样的多项式不可能在所有仿射有理点上消失，这可按变量个数用一元多项式的根数界归纳。又 $P^\vee$ 不可约，包含泛点的可构造子集含一个非空开集。故该步骤先由可构造性得开集，再用无限域取得有理点，两步不能混同。

在 $s$ 附近平凡化参数射影丛，并选取含 $\xi_0$ 的标准仿射图。将 $\xi_0$ 在这张图的坐标分别从 $k(s)$ 提升到 $\mathcal O_{S,s}$；有限个局部环元素可在同一个开邻域上表示成正则函数，于是得到 $\xi$。当 $E$ 开时，把邻域限制到 $\xi^{-1}(E)$ 即可。仅知道 $E$ 可构造，不能保证这个逆像含邻域。

\paragraph{有限域与有限 \'{e}tale 覆盖.}
例如 $\mathbf P^1_{\mathbf F_q}\setminus\mathbf P^1(\mathbf F_q)$ 非空，却没有 $\mathbf F_q$-有理点。非空开集中的闭点具有有限剩余域扩张，而有限域是完美的。若 $k'=k(s)[T]/(\overline F)$，其中 $\overline F$ 首一、可分且不可约，在 $s$ 的一个仿射邻域 $\Spec A$ 上把系数提升为首一多项式 $F\in A[T]$。$A[T]/(F)$ 是正秩有限自由 $A$-模；其判别式在 $s$ 不为零，缩小基底使判别式可逆，就得到所需有限 étale 满射，且在 $s$ 的纤维正是 $\Spec k'$。所选参数经此换底成为 $k'$-有理点。这里仍要在新基底上核验好参数开性及平坦性判据，不从这个换底本身推出任意性质 $\calP$ 的保持。可对照 \href{https://stacks.math.columbia.edu/tag/02LF}{Stacks Project，有限可分剩余域的 étale 邻域}。

\clearpage
\section*{法文划销旧稿：不并入正文}
\paragraph{法26上部：关于第4节的改写建议.}
作者说，旧4.6(b)、(c)的结论已包含于旧4.7(b)、(c)。起初因后者的证明似乎需要更细致的结果（用15.2.2的形式），才认为从证明角度可以保留重复。若只想证明相应集合 $U$（分别为 $E$）开，而不证明 $\GG|_U$ 相对 $P^\vee$ 平坦，可以改走第12段使用的赋值判据。作者建议删去4.6(b)、(c)，将4.7(b)、(c)仅表述成开性（同时删4.7(a)），再添命题或推论单列平坦性并用15.2.2。行间可见手记“4.8”。全段被划销，编号均按废稿记，不据此改写现行第4、5节。正式 EGA~IV 第三分册（15.2.2）确是带附加条件的平坦性定理，见\href{https://www.numdam.org/item/PMIHES_1966__28__5_0/}{印刷页226}；不是“开便平坦”。\footnote{废稿局部断字，以上以能辨认的三条编辑建议及当时编号为准。正式（15.2.2）要求支撑、相应泛点处的普遍开性、模纤维几何约化和底局部环约化等条件；未宣称废稿对这些条件的应用已经完证。}

\paragraph{法27下部：与旧4.6、4.7相关的另一则剪裁稿.}
“下面给出……例子”一句未删。其后被划销的片段准备讨论如何放宽涉及超平面和 $\Supp\FF_s$ 的几何不可约分支的假设，拟给 $P^\vee$ 取一个开子集 $W$ 后改写条件。右端和末句不全，本批只登记可辨方向，不据片段补造完整假设。法28起笔是另一个未删的关于4.2与4.3的补记，留待下批。
\end{document}

